\chapter{Functional Analysis}\label{functional-analysis}
\chapterstatus{complete}

\section{Introduction}\label{functional-analysis:section-introduction}

Hodge theory rests on an analytic theorem: on a compact Riemannian manifold every cohomology class has a unique harmonic
representative. The proof solves the equation $\Delta u = f$ in a space of generalised functions and then shows that the solution is
smooth. The functional analysis needed for the first step is in this chapter: Hilbert spaces and orthogonal projections, the
consequences of Baire's theorem, compact operators, the spectral theorem for compact self-adjoint operators
(Theorem~\ref{functional-analysis:theorem-spectral-compact}) and the Fredholm alternative
(Theorem~\ref{functional-analysis:theorem-fredholm-alternative}). The guiding theme is what survives from finite-dimensional linear algebra and what does not: in infinite
dimensions linear maps need not be continuous, norms need not be equivalent, closed bounded sets need not be compact, injective operators need not be surjective, and bases have to be replaced by
orthonormal bases, whose expansions are convergent series. We also introduce Fr\'echet spaces, such as the space of smooth functions, whose topology is given by countably many seminorms; they are the
natural home of distributions (Chapter~\ref{distributions}). References: \cite{brezis-functional}, \cite{rudin-fa}, \cite{reed-simon-1}.

\noindent\textbf{Prerequisites.} Chapter~\ref{measure-theory} (Measure and Integration) and Chapter~\ref{linear-algebra} (Linear Algebra).

All vector spaces are over $\mathbb{K} = \RR$ or $\CC$.

\section{Banach spaces and bounded operators}\label{functional-analysis:section-banach}

\begin{definition}\label{functional-analysis:definition-bounded-operator}
A linear map $T \colon X \to Y$ between normed spaces is \emph{bounded} if $\|T\| = \sup_{\|x\| \leq 1}\|Tx\| < \infty$; equivalently (as for finite dimensions,
Lemma~\ref{real-analysis:lemma-norms}) $T$ is continuous. The bounded operators form a normed space $B(X, Y)$, and $X^* = B(X, \mathbb{K})$ is the \emph{dual space}.
\end{definition}

\begin{proposition}\label{functional-analysis:proposition-banach-examples}
\begin{enumerate}
\item If $Y$ is complete, so is $B(X, Y)$; in particular $X^*$ is always a Banach space.
\item $C(K, \mathbb{K})$ with the supremum norm ($K$ compact Hausdorff), $L^p(\mu)$ for $1 \leq p \leq \infty$, and the sequence spaces $\ell^p$ are Banach spaces.
\item (Riesz) If the closed unit ball of a normed space $X$ is compact, then $X$ is finite-dimensional.
\end{enumerate}
\end{proposition}

\begin{proof}
(1) A Cauchy sequence $(T_n)$ is pointwise Cauchy, and the pointwise limit $T$ is linear with $\|T - T_n\| \leq \limsup_m\|T_m - T_n\| \to 0$. (2)
Proposition~\ref{general-topology:proposition-complete-examples} and Theorem~\ref{measure-theory:theorem-riesz-fischer}. (3) Suppose $X$ is infinite-dimensional. We construct unit vectors $x_n$ with
$\|x_n - x_m\| \geq 1/2$: given the finite-dimensional, hence closed (all norms on it are equivalent and it is complete), subspace $Y_n$ spanned by $x_1, \ldots, x_n$, choose $x \notin Y_n$ and $y \in Y_n$ with
$\|x - y\| \leq 2\operatorname{dist}(x, Y_n)$; then $x_{n+1} = (x - y)/\|x - y\|$ has distance at least $1/2$ from $Y_n$. Such a sequence has no convergent subsequence.
\end{proof}

\begin{example}\label{functional-analysis:example-norms}
\begin{enumerate}
\item On $\mathbb{K}^n$ the norms $\|x\|_p = (\sum_i|x_i|^p)^{1/p}$ ($1 \leq p < \infty$) and $\|x\|_\infty = \max_i|x_i|$ are all equivalent (Lemma~\ref{real-analysis:lemma-norms}); for instance
$\|x\|_\infty \leq \|x\|_2 \leq \|x\|_1 \leq n\|x\|_\infty$. Their unit balls are a cube, a ball and a cross-polytope.
\item On $C([0, 1])$ the norms $\|f\|_\infty$ and $\|f\|_1 = \int_0^1|f|$ are not equivalent: $\|f\|_1 \leq \|f\|_\infty$, but $f_n(x) = x^n$ has $\|f_n\|_\infty = 1$ and $\|f_n\|_1 = 1/(n + 1)$. So
$f_n \to 0$ in one norm and not in the other, and $C([0, 1])$ is complete for $\|\cdot\|_\infty$ but not for $\|\cdot\|_1$ (Exercise~\ref{functional-analysis:exercise-l1-incomplete}).
\item \emph{Norms of operators.} A matrix $A \colon \mathbb{K}^n \to \mathbb{K}^m$ with Euclidean norms has $\|A\| = $ its largest singular value, the square root of the largest eigenvalue of $A^*A$
(Proposition~\ref{linear-algebra:proposition-svd}). For $g \in L^\infty([0, 1])$ the multiplication operator $M_gf = gf$ on $L^2([0, 1])$ has $\|M_g\| = \|g\|_\infty$
(Exercise~\ref{functional-analysis:exercise-multiplication-operator}). For a continuous kernel $k$ on $[0, 1]^2$, $(Tf)(x) = \int_0^1k(x, y)f(y)\,dy$ is bounded on $C([0, 1])$ with
$\|T\| \leq \sup_x\int_0^1|k(x, y)|\,dy$.
\end{enumerate}
\end{example}

\begin{lemma}\label{functional-analysis:lemma-banach-series}
A normed space $X$ is complete if and only if every absolutely convergent series ($\sum\|x_n\| < \infty$) converges. If $X$ is a Banach space and $M \subseteq X$ a closed subspace, the quotient $X/M$ with
$\|x + M\| = \operatorname{dist}(x, M)$ is a Banach space, and the projection $X \to X/M$ has norm at most $1$.
\end{lemma}

\begin{proof}
If $X$ is complete, the partial sums of an absolutely convergent series are Cauchy (Proposition~\ref{real-analysis:proposition-series}(1)). Conversely, given a Cauchy sequence choose $n_1 < n_2 < \cdots$ with
$\|x_{n_{k+1}} - x_{n_k}\| < 2^{-k}$; the series $x_{n_1} + \sum_k(x_{n_{k+1}} - x_{n_k})$ converges absolutely, so the subsequence converges, and then so does the Cauchy sequence. For the quotient,
$\|x + M\| = 0$ forces $x \in \overline{M} = M$, and the other axioms are clear, so it is a norm. If $\sum\|x_n + M\| < \infty$, choose $m_n \in M$ with $\|x_n - m_n\| \leq \|x_n + M\| + 2^{-n}$; then
$\sum(x_n - m_n)$ converges absolutely in $X$, to some $x$, and $\sum(x_n + M)$ converges to $x + M$.
\end{proof}

\begin{remark}\label{functional-analysis:remark-separable-bases}
A normed space is \emph{separable} if it has a countable dense subset. $\ell^p$ and $L^p([0, 1])$ for $p < \infty$, and $C([0, 1])$, are separable (finitely supported sequences with rational entries;
step functions or polynomials with rational coefficients, Corollary~\ref{real-analysis:corollary-weierstrass}). $\ell^\infty$ is not: the indicator sequences $\mathbf{1}_S$ of the uncountably many
subsets $S \subseteq \NN$ are at mutual distance $1$. A \emph{Hamel basis}, a basis in the sense of linear algebra, of an infinite-dimensional Banach space is uncountable: if $e_1, e_2, \ldots$ were one,
$X$ would be the countable union of the closed subspaces $\operatorname{span}(e_1, \ldots, e_n)$, which have empty interior, contradicting Baire's theorem
(Theorem~\ref{general-topology:theorem-baire}). This is why in analysis one works with expansions $x = \sum_n c_ne_n$ as convergent series instead: orthonormal bases in Hilbert spaces
(Theorem~\ref{functional-analysis:theorem-orthonormal-basis}) and, more generally, \emph{Schauder bases}.
\end{remark}

\section{Hilbert spaces}\label{functional-analysis:section-hilbert}

\begin{definition}\label{functional-analysis:definition-hilbert}
A \emph{Hilbert space} is a vector space with an inner product (Definition~\ref{linear-algebra:definition-inner-product}) that is complete for the norm $\|x\| = \langle x, x \rangle^{1/2}$. Vectors are
\emph{orthogonal} if $\langle x, y \rangle = 0$; for $S \subseteq H$, $S^\perp = \{x : \langle x, s \rangle = 0 \ \forall s \in S\}$, a closed subspace.
\end{definition}

\begin{example}\label{functional-analysis:example-hilbert}
$L^2(\mu)$ with $\langle f, g \rangle = \int f\bar g$ (Theorem~\ref{measure-theory:theorem-riesz-fischer}); $\ell^2$; finite-dimensional inner product spaces. The space $C([0,1])$ with the $L^2$ inner product is
not complete, and $L^2$ is its completion (Exercise~\ref{general-topology:exercise-completion}).
\end{example}

\begin{theorem}\label{functional-analysis:theorem-projection}
Let $H$ be a Hilbert space and $C \subseteq H$ nonempty, closed and convex. Every $x \in H$ has a unique nearest point $P_Cx \in C$. If $M$ is a closed subspace, $P_M$ is linear with $\|P_M\| \leq 1$,
$x - P_Mx \in M^\perp$, and $H = M \oplus M^\perp$. Moreover $(M^\perp)^\perp = M$, and $M^\perp = 0$ if and only if $M$ is dense.
\end{theorem}

\begin{proof}
Let $d = \inf_{c \in C}\|x - c\|$ and $c_k \in C$ with $\|x - c_k\| \to d$. By the parallelogram law,
$\|c_k - c_l\|^2 = 2\|x - c_k\|^2 + 2\|x - c_l\|^2 - 4\|x - \tfrac{c_k + c_l}{2}\|^2 \leq 2\|x - c_k\|^2 + 2\|x - c_l\|^2 - 4d^2 \to 0$, using $\tfrac{c_k+c_l}{2} \in C$. So $(c_k)$ is Cauchy and converges to some
$c \in C$ with $\|x - c\| = d$; uniqueness follows from the same identity. For a closed subspace $M$ and $m \in M$, $t \in \mathbb{K}$, the inequality $\|x - P_Mx - tm\| \geq \|x - P_Mx\|$ forces
$\langle x - P_Mx, m \rangle = 0$. Hence $x = P_Mx + (x - P_Mx) \in M + M^\perp$, and $M \cap M^\perp = 0$, which also gives linearity and $\|P_M\| \leq 1$ (Pythagoras). Finally $M \subseteq (M^\perp)^\perp$, and
if $x \in (M^\perp)^\perp$ then $x - P_Mx \in M^\perp \cap (M^\perp)^\perp = 0$. If $M$ is a subspace, $\overline{M}^\perp = M^\perp$, so $M^\perp = 0$ if and only if $\overline{M} = H$.
\end{proof}

\begin{theorem}[Riesz representation]\label{functional-analysis:theorem-riesz-representation}
For every $\Lambda \in H^*$ there is a unique $y \in H$ with $\Lambda(x) = \langle x, y \rangle$ for all $x$, and $\|\Lambda\| = \|y\|$. Consequently every bounded operator $T \colon H \to H$ has a unique
\emph{adjoint} $T^*$ with $\langle Tx, y \rangle = \langle x, T^*y \rangle$; it satisfies $\|T^*\| = \|T\|$ and $(\ker T)^\perp = \overline{\operatorname{im}T^*}$.
\end{theorem}

\begin{proof}
Existence and uniqueness of $y$ are Lemma~\ref{measure-theory:lemma-hilbert-riesz}, whose proof used only completeness; the norm equality follows from Cauchy--Schwarz and $\Lambda(y) = \|y\|^2$. For fixed
$y$, $x \mapsto \langle Tx, y \rangle$ is a bounded functional, so it is $\langle x, T^*y \rangle$ for a unique $T^*y$, and $T^*$ is linear and bounded with $\|T^*\| = \|T\|$ (from
$\|T\| = \sup_{\|x\|, \|y\| \leq 1}|\langle Tx, y \rangle|$, which is symmetric in $T$ and $T^*$). Finally $x \in \ker T$ if and only if $\langle x, T^*y \rangle = 0$ for all $y$, that is, $x \in (\operatorname{im}T^*)^\perp$; take orthogonal
complements.
\end{proof}

\begin{proposition}[Bessel's inequality]\label{functional-analysis:proposition-bessel}
Let $(e_i)_{i \in I}$ be an orthonormal family in an inner product space $H$. For every $x \in H$,
\[
\sum_{i \in I}|\langle x, e_i \rangle|^2 \leq \|x\|^2;
\]
in particular only countably many coefficients $\langle x, e_i\rangle$ are nonzero, and $\langle x, e_n\rangle \to 0$ for an orthonormal sequence.
\end{proposition}

\begin{proof}
For a finite $J \subseteq I$ let $y = \sum_{i \in J}\langle x, e_i \rangle e_i$. Then $\langle x - y, e_j\rangle = 0$ for $j \in J$, so $x - y \perp y$ and, by Pythagoras,
$\|x\|^2 = \|y\|^2 + \|x - y\|^2 \geq \|y\|^2 = \sum_{i \in J}|\langle x, e_i \rangle|^2$. A sum of nonnegative terms is the supremum of its finite partial sums. For each $k$ at most $k\|x\|^2$ indices
have $|\langle x, e_i\rangle|^2 \geq 1/k$, so the set of nonzero coefficients is a countable union of finite sets.
\end{proof}

The vector $y$ in the proof is the best approximation of $x$ by combinations of $e_i$, $i \in J$: $\|x - y\|^2 = \|x\|^2 - \sum_{i \in J}|\langle x, e_i\rangle|^2$, and equality in Bessel's inequality for all $x$ is
Parseval's identity, which characterises orthonormal bases.

\begin{theorem}[Orthonormal bases]\label{functional-analysis:theorem-orthonormal-basis}
Let $(e_i)_{i \in I}$ be an orthonormal family in a Hilbert space $H$. The following are equivalent: its span is dense; $x = \sum_i\langle x, e_i \rangle e_i$ for all $x$ (with at most countably many nonzero
terms, in any order); $\|x\|^2 = \sum_i|\langle x, e_i \rangle|^2$ for all $x$ (Parseval). Such a family is an \emph{orthonormal basis}; every Hilbert space has one, and $H$ is separable if and only if it has a
countable one, in which case $H \cong \ell^2(I)$.
\end{theorem}

\begin{proof}
For finite $J \subseteq I$, $P_Jx = \sum_{i \in J}\langle x, e_i \rangle e_i$ is the orthogonal projection onto the span of $(e_i)_{i \in J}$, so $\sum_{i \in J}|\langle x, e_i \rangle|^2 \leq \|x\|^2$ (Bessel);
hence at most countably many coefficients are nonzero and, by completeness, $\sum_i\langle x, e_i \rangle e_i$ converges (its partial sums are Cauchy by Pythagoras), to the projection of $x$ onto the closed
span $M$. So the three conditions each say that $M = H$, using $\|x\|^2 = \|P_Mx\|^2 + \|x - P_Mx\|^2$. Zorn's lemma gives a maximal orthonormal family, whose closed span $M$ satisfies $M^\perp = 0$ (a unit
vector in $M^\perp$ would enlarge the family), so $M = H$. If $H$ is separable, an orthonormal family is discrete ($\|e_i - e_j\| = \sqrt2$) hence countable. The map $x \mapsto (\langle x, e_i \rangle)$ is
then an isometric isomorphism onto $\ell^2(I)$.
\end{proof}

\begin{example}\label{functional-analysis:example-fourier}
The functions $e_n(t) = e^{2\pi int}$, $n \in \ZZ$, form an orthonormal basis of $L^2(S^1)$, where $S^1 = \RR/\ZZ$ with Lebesgue measure: they are orthonormal by direct computation, their span (the
trigonometric polynomials) is dense in $C(S^1)$ for the supremum norm (Corollary~\ref{real-analysis:corollary-weierstrass}), hence in $L^2$, and $C(S^1)$ is dense in $L^2(S^1)$: by
Proposition~\ref{measure-theory:proposition-density} the continuous functions with compact support in $(0, 1)$ are dense in $L^2([0, 1])$, and they extend to continuous functions on $S^1$. So every $f \in L^2(S^1)$ is the $L^2$-limit of its Fourier series $\sum_n\hat f(n)e_n$ with $\hat f(n) = \int_0^1f(t)e^{-2\pi int}\,dt$, and
$\|f\|_2^2 = \sum_n|\hat f(n)|^2$.
\end{example}

\begin{example}\label{functional-analysis:example-bases-l2}
\emph{Orthonormal bases of $L^2([0, 1])$.} Up to the null set $\{1\}$, $[0, 1)$ is the circle $\RR/\ZZ$, so Example~\ref{functional-analysis:example-fourier} gives the basis $e^{2\pi int}$, $n \in \ZZ$.
\begin{enumerate}
\item \emph{Real Fourier basis.} Taking real and imaginary parts, $1$, $\sqrt2\cos(2\pi nt)$, $\sqrt2\sin(2\pi nt)$ ($n \geq 1$) is an orthonormal basis of the real space $L^2([0, 1])$: these functions are
orthonormal ($\int_0^1\cos^2(2\pi nt)\,dt = \frac12$), and they span the same space of trigonometric polynomials, over $\CC$, as the $e^{\pm2\pi int}$.
\item \emph{Legendre polynomials.} Gram--Schmidt (Theorem~\ref{linear-algebra:theorem-gram-schmidt}) applied to $1, x, x^2, \ldots$ gives orthonormal polynomials $P_0, P_1, P_2, \ldots$ with
$\deg P_n = n$. With $\langle f, g\rangle = \int_0^1fg$: $P_0 = 1$; $x - \langle x, 1\rangle = x - \frac12$ has norm$^2$ $\int_0^1(x - \frac12)^2 = \frac1{12}$, so $P_1 = \sqrt3(2x - 1)$; and
$x^2 - \langle x^2, P_0\rangle P_0 - \langle x^2, P_1\rangle P_1 = x^2 - \frac13 - \frac{\sqrt3}{6}\sqrt3(2x - 1) = x^2 - x + \frac16$, whose norm$^2$ is $\frac1{180}$, so $P_2 = \sqrt5(6x^2 - 6x + 1)$.
They form an orthonormal basis: their span is the space of all polynomials, which is dense in $C([0, 1])$ for the supremum norm (Corollary~\ref{real-analysis:corollary-weierstrass}), hence in
$L^2$ (since $\|f\|_2 \leq \|f\|_\infty$ and $C([0, 1])$ is dense in $L^2$, Proposition~\ref{measure-theory:proposition-density}), and Theorem~\ref{functional-analysis:theorem-orthonormal-basis} applies.
\item \emph{The Haar basis.} Let $h = \mathbf{1}_{[0, 1/2)} - \mathbf{1}_{[1/2, 1)}$ and $h_{n,k}(t) = 2^{n/2}h(2^nt - k)$ for $n \geq 0$, $0 \leq k < 2^n$. Together with $1$ these are orthonormal: two of them
either have disjoint supports, or the support of one lies in an interval where the other is constant, and then the integral of the first vanishes. The $2^N$ functions $1$ and $h_{n,k}$ with $n < N$ lie
in the $2^N$-dimensional space of functions constant on the dyadic intervals $[j2^{-N}, (j + 1)2^{-N})$ and are independent, so they span it. Such step functions approximate continuous functions
uniformly, so the span is dense, and this is an orthonormal basis; unlike the Fourier basis it consists of functions with small support, the starting point of wavelets.
\end{enumerate}
\end{example}

\begin{example}\label{functional-analysis:example-basel}
\emph{Parseval and $\sum 1/n^2$.} For $f(t) = t$ on $[0, 1)$, $\hat f(0) = \frac12$ and, integrating by parts, $\hat f(n) = \int_0^1te^{-2\pi int}\,dt = \bigl[\frac{te^{-2\pi int}}{-2\pi in}\bigr]_0^1 = \frac{i}{2\pi n}$
for $n \neq 0$. Parseval's identity gives $\frac13 = \int_0^1t^2\,dt = \frac14 + \sum_{n \neq 0}\frac{1}{4\pi^2n^2}$, that is, $\sum_{n \geq 1}\frac1{n^2} = \frac{\pi^2}6$. Note that the Fourier series of $f$
converges to $f$ in $L^2$, not pointwise at $t = 0$, where the $1$-periodic extension of $f$ jumps.
\end{example}

\section{The Hahn--Banach and Baire theorems}\label{functional-analysis:section-hahn-banach}

\begin{theorem}[Hahn--Banach]\label{functional-analysis:theorem-hahn-banach}
Let $X$ be a vector space over $\mathbb{K}$ and $p \colon X \to \RR$ with $p(x + y) \leq p(x) + p(y)$ and $p(\lambda x) = |\lambda|p(x)$ (a seminorm). Every linear functional $\ell$ on a subspace $Y$ with
$|\ell| \leq p$ on $Y$ extends to a linear functional on $X$ with $|\ell| \leq p$. Consequently, for a normed $X$ and $x_0 \neq 0$ there is $\Lambda \in X^*$ with $\|\Lambda\| = 1$ and $\Lambda(x_0) = \|x_0\|$;
in particular $X^*$ separates points and $\|x\| = \sup_{\|\Lambda\| \leq 1}|\Lambda(x)|$.
\end{theorem}

\begin{proof}
Real case: consider pairs $(Z, m)$ of a subspace $Z \supseteq Y$ and an extension $m \leq p$, ordered by extension; chains have upper bounds, so Zorn's lemma gives a maximal $(Z, m)$. If $x \notin Z$, then for
$z, z' \in Z$, $m(z) + m(z') = m(z + z') \leq p(z + x) + p(z' - x)$, so $\sup_{z'}(m(z') - p(z' - x)) \leq \inf_z(p(z + x) - m(z))$; choosing $c$ between them extends $m$ to $Z + \mathbb{K}x$ by $m(z + tx) = m(z) + tc$,
contradicting maximality. Complex case: apply the real case to $\operatorname{Re}\ell$ and set $\ell(x) = (\operatorname{Re}\ell)(x) - i(\operatorname{Re}\ell)(ix)$. For the consequence take
$Y = \mathbb{K}x_0$, $\ell(tx_0) = t\|x_0\|$ and $p = \|\cdot\|$.
\end{proof}

\begin{theorem}\label{functional-analysis:theorem-baire-consequences}
Let $X$, $Y$ be Banach spaces.
\begin{enumerate}
\item (Uniform boundedness) If $\mathcal{T} \subseteq B(X, Y)$ satisfies $\sup_{T \in \mathcal{T}}\|Tx\| < \infty$ for every $x$, then $\sup_{T \in \mathcal{T}}\|T\| < \infty$.
\item (Open mapping) A surjective $T \in B(X, Y)$ is open; if it is bijective, $T^{-1}$ is bounded.
\item (Closed graph) A linear $T \colon X \to Y$ whose graph is closed in $X \times Y$ is bounded.
\end{enumerate}
\end{theorem}

\begin{proof}
(1) The sets $F_n = \{x : \|Tx\| \leq n \ \forall T \in \mathcal{T}\}$ are closed with union $X$, so by Baire's theorem (Theorem~\ref{general-topology:theorem-baire}) some $F_n$ contains a ball
$\overline{B}(x_0, r)$; then $\|Tx\| \leq 2n/r$ for $\|x\| \leq 1$, using $x = ((x_0 + rx) - x_0)/r$.

(2) By Baire's theorem applied to $Y = \bigcup_n\overline{T(B(0, n))}$, the closure of $T(B(0, 1))$ contains a ball $B(y_0, \delta)$; by symmetry and convexity it contains $B(0, \delta)$. A standard iteration
removes the closure: given $\|y\| < \delta$, choose $x_1$ with $\|x_1\| < 1$ and $\|y - Tx_1\| < \delta/2$, then $x_2$ with $\|x_2\| < 1/2$ and $\|y - Tx_1 - Tx_2\| < \delta/4$, and so on; the series
$\sum x_k$ converges to some $x$ with $\|x\| < 2$ and $Tx = y$. So $T(B(0, 2)) \supseteq B(0, \delta)$, which gives openness by translation and scaling.

(3) The graph $G$ is a Banach space, and the projection $G \to X$ is bounded and bijective, so its inverse $x \mapsto (x, Tx)$ is bounded by (2).
\end{proof}

\begin{counterexample}\label{functional-analysis:counterexample-unbounded}
\begin{enumerate}
\item \emph{An unbounded operator.} Differentiation $D \colon C^1([0, 1]) \to C([0, 1])$, both with the supremum norm, is linear but not bounded: $\|x^n\|_\infty = 1$ and $\|D(x^n)\|_\infty = n$. It is closed (if $f_k \to f$ and $f_k' \to g$
uniformly, then $f' = g$, Theorem~\ref{real-analysis:theorem-uniform-derivative}), which does not contradict the closed graph theorem because $C^1([0, 1])$ is not complete for the supremum norm.
\item \emph{An incomplete normed space.} The polynomials, with the supremum norm on $[0, 1]$: the partial sums of $\sum x^k/k!$ form a Cauchy sequence whose uniform limit $e^x$ is not a polynomial, since all its derivatives are nonzero.
\item \emph{Completeness is needed for uniform boundedness.} On the space $c_{00}$ of finitely supported sequences with the supremum norm, the functionals $T_n(x) = nx_n$ satisfy $\sup_n|T_nx| < \infty$ for each $x$, but $\|T_n\| = n$.
\item \emph{Discontinuous linear functionals.} Every infinite-dimensional normed space $X$ has one. Choose linearly independent unit vectors $e_1, e_2, \ldots$ and extend them, by Zorn's lemma
(Theorem~\ref{sets:theorem-choice-equivalents}), to a maximal linearly independent set $B$, which is a basis of $X$ as a vector space. The linear functional with $\ell(e_n) = n$ and $\ell = 0$ on $B \setminus \{e_1, e_2, \ldots\}$ is unbounded.
In finite dimensions every linear functional is continuous (Lemma~\ref{real-analysis:lemma-norms}).
\end{enumerate}
\end{counterexample}

\section{Compact operators}\label{functional-analysis:section-compact-operators}

\begin{definition}\label{functional-analysis:definition-compact-operator}
$T \in B(X, Y)$ is \emph{compact} if the closure of $T(B(0,1))$ is compact; equivalently, every bounded sequence $(x_n)$ has a subsequence with $(Tx_{n_k})$ convergent.
\end{definition}

\begin{proposition}\label{functional-analysis:proposition-compact-operators}
\begin{enumerate}
\item Compact operators form a closed subspace of $B(X, Y)$, and $STR$ is compact if $T$ is compact and $S$, $R$ are bounded.
\item Operators of finite rank are compact; on Hilbert spaces, $T$ is compact if and only if it is a norm limit of finite rank operators.
\item If $H$ is a Hilbert space, $T \in B(H)$ is compact if and only if $T^*$ is.
\end{enumerate}
\end{proposition}

\begin{proof}
(1) Sums and scalar multiples are clear, and compositions with bounded operators map bounded sets to bounded sets. If $T_n \to T$ with $T_n$ compact, then $T(B)$ is totally bounded: given $\varepsilon$, choose
$n$ with $\|T - T_n\| < \varepsilon$ and cover $T_n(B)$ by finitely many $\varepsilon$-balls; so $\overline{T(B)}$ is compact (Theorem~\ref{general-topology:theorem-compact-metric}). (2) A bounded set in a
finite-dimensional space is relatively compact (Theorem~\ref{general-topology:theorem-heine-borel}). Conversely let $T$ be compact and $\varepsilon > 0$; cover $\overline{T(B)}$ by finitely many balls
$B(y_i, \varepsilon)$, let $P$ be the orthogonal projection onto the span of the $y_i$ (Theorem~\ref{functional-analysis:theorem-projection}); then $\|Tx - PTx\| \leq \|Tx - y_i\| < \varepsilon$ for a suitable
$i$, since $P$ is a projection and $\|Tx - PTx\| = \operatorname{dist}(Tx, \operatorname{im}P)$. So $\|T - PT\| \leq \varepsilon$ with $PT$ of finite rank. (3) follows from (2), since $(PT)^* = T^*P$ has finite
rank and $\|T^* - (PT)^*\| = \|T - PT\|$.
\end{proof}

\begin{theorem}[Spectral theorem for compact self-adjoint operators]\label{functional-analysis:theorem-spectral-compact}
Let $H \neq 0$ be a Hilbert space and $T \in B(H)$ compact and self-adjoint ($T = T^*$).
\begin{enumerate}
\item $\|T\| = \sup_{\|x\| = 1}|\langle Tx, x \rangle|$, and $\|T\|$ or $-\|T\|$ is an eigenvalue of $T$.
\item $H$ has an orthonormal basis consisting of eigenvectors of $T$. The nonzero eigenvalues are real, form a finite or countable set with $0$ as the only possible accumulation point, and each nonzero
eigenvalue has finite-dimensional eigenspace.
\end{enumerate}
\end{theorem}

\begin{proof}
(1) Let $s = \sup_{\|x\|=1}|\langle Tx, x \rangle| \leq \|T\|$. For $x, y$ with $\|x\|, \|y\| \leq 1$, the identity
$\langle T(x + y), x + y \rangle - \langle T(x - y), x - y \rangle = 4\operatorname{Re}\langle Tx, y \rangle$ (using self-adjointness) and the parallelogram law give
$4\operatorname{Re}\langle Tx, y \rangle \leq s(\|x + y\|^2 + \|x - y\|^2) = 2s(\|x\|^2 + \|y\|^2) \leq 4s$, so $\|T\| = \sup|\langle Tx, y \rangle| \leq s$. Choose unit vectors $x_n$ with
$\langle Tx_n, x_n \rangle \to \lambda$ where $|\lambda| = \|T\|$ (note $\langle Tx, x \rangle$ is real). Then
$\|Tx_n - \lambda x_n\|^2 = \|Tx_n\|^2 - 2\lambda\langle Tx_n, x_n \rangle + \lambda^2 \leq 2\lambda^2 - 2\lambda\langle Tx_n, x_n \rangle \to 0$. By compactness, after passing to a subsequence
$Tx_n \to y$; if $\lambda \neq 0$ then $x_n \to y/\lambda$, so $y/\lambda$ is a unit vector with $Ty/\lambda = \lambda y/\lambda$. If $\lambda = 0$ then $T = 0$ and every vector is an eigenvector.

(2) Eigenvalues are real and eigenvectors for distinct eigenvalues are orthogonal, as in finite dimensions. If $|\lambda| \geq \varepsilon > 0$, an orthonormal family in the $\lambda$-eigenspace is mapped by $T$ to
a family of vectors at mutual distance $|\lambda|\sqrt2$, so it must be finite by compactness; the same argument shows that only finitely many eigenvalues satisfy $|\lambda| \geq \varepsilon$. Let $M$ be the
closed span of all eigenvectors. It is $T$-invariant, hence so is $M^\perp$ (if $x \in M^\perp$ and $m \in M$, $\langle Tx, m \rangle = \langle x, Tm \rangle = 0$), and $T|_{M^\perp}$ is compact and
self-adjoint; by (1) it has an eigenvector unless $M^\perp = 0$ or $T|_{M^\perp} = 0$. An eigenvector in $M^\perp$ is impossible, and if $T|_{M^\perp} = 0$ then $M^\perp$ consists of eigenvectors for the
eigenvalue $0$, so $M^\perp \subseteq M$ and $M^\perp = 0$. Hence $M = H$, and an orthonormal basis of each eigenspace, together, is an orthonormal basis of $H$ by
Theorem~\ref{functional-analysis:theorem-orthonormal-basis}.
\end{proof}

\begin{theorem}[Fredholm alternative]\label{functional-analysis:theorem-fredholm-alternative}
Let $H$ be a Hilbert space, $K \in B(H)$ compact and $T = I - K$.
\begin{enumerate}
\item $\ker T$ is finite-dimensional and $\operatorname{im}T$ is closed, with $\operatorname{im}T = (\ker T^*)^\perp$.
\item $T$ is injective if and only if it is surjective.
\item $\dim\ker T = \dim\ker T^*$; so the equation $Tx = y$ is solvable if and only if $y \perp \ker T^*$, and the solution is unique up to $\ker T$, which has the same finite dimension as $\ker T^*$.
\end{enumerate}
\end{theorem}

\begin{proof}
(1) On $\ker T$ we have $K = I$, so the unit ball of $\ker T$ is compact and $\ker T$ is finite-dimensional (Proposition~\ref{functional-analysis:proposition-banach-examples}(3)). For closedness, let
$M = (\ker T)^\perp$; it suffices to show $\|Tx\| \geq c\|x\|$ on $M$ for some $c > 0$, since then images of Cauchy sequences are Cauchy. If not, there are unit vectors $x_n \in M$ with $Tx_n \to 0$; passing to a
subsequence, $Kx_n \to y$, so $x_n = Tx_n + Kx_n \to y$ with $\|y\| = 1$, $y \in M$ and $Ty = 0$, a contradiction. Then $\operatorname{im}T$ is closed and
$\operatorname{im}T = \overline{\operatorname{im}T} = (\ker T^*)^\perp$ by Theorem~\ref{functional-analysis:theorem-riesz-representation} applied to $T^*$.

(2) Suppose $T$ is injective but not surjective. Then $H \supsetneq TH \supsetneq T^2H \supsetneq \cdots$ is a strictly decreasing chain of closed subspaces (each $T^kH$ is closed by (1) applied to
$I - K_k$ with $K_k = I - T^k$ compact, since $T^k = (I - K)^k = I - K_k$ with $K_k$ a sum of products involving $K$). Choose unit vectors $u_k \in T^kH \cap (T^{k+1}H)^\perp$. For $l > k$,
$Ku_k - Ku_l = u_k - (Tu_k + u_l - Tu_l)$ and $Tu_k + u_l - Tu_l \in T^{k+1}H$, so $\|Ku_k - Ku_l\| \geq 1$, contradicting compactness. Conversely, if $T$ is surjective, then $T^*$ is injective: indeed
$\ker T^* = (\operatorname{im}T)^\perp = 0$; applying the previous argument to $T^* = I - K^*$ (with $K^*$ compact) shows $T^*$ is surjective, hence $\ker T = (\operatorname{im}T^*)^\perp = 0$.

(3) This is the statement that the index of $I - K$ is zero; see \cite[Theorem 6.6]{brezis-functional}.
\end{proof}

\begin{remark}\label{functional-analysis:remark-fredholm-operators}
$T \in B(X, Y)$ is \emph{Fredholm} if $\ker T$ and $\operatorname{coker}T = Y/\operatorname{im}T$ are finite-dimensional; then $\operatorname{im}T$ is closed, and the \emph{index}
$\operatorname{ind}T = \dim\ker T - \dim\operatorname{coker}T$ is unchanged by compact perturbations and is locally constant in $T$ \cite[Chapter 6]{brezis-functional}. Elliptic operators on compact
manifolds are Fredholm (Chapter~\ref{elliptic-operators}); the index is then computed by the Atiyah--Singer theorem, and its vanishing for the Laplacian is what produces harmonic representatives.
\end{remark}

\begin{counterexample}\label{functional-analysis:counterexample-no-eigenvalues}
Self-adjointness is needed in the spectral theorem.
\begin{enumerate}
\item The \emph{Volterra operator} $Vf(x) = \int_0^xf(t)\,dt$ on $L^2([0, 1])$ is compact but has no eigenvalues. It is compact: by the Cauchy--Schwarz inequality $|Vf(x) - Vf(y)| \leq |x - y|^{1/2}\|f\|_2$ and $|Vf(x)| \leq \|f\|_2$, so
$V$ maps the unit ball to a bounded equicontinuous set of continuous functions, which is relatively compact in $C([0, 1])$ by the Arzel\`a--Ascoli theorem (Theorem~\ref{real-analysis:theorem-arzela-ascoli}), hence in $L^2$. If
$Vf = \lambda f$ with $\lambda \neq 0$, then $f = \lambda^{-1}Vf$ is continuous, hence continuously differentiable, with $f' = \lambda^{-1}f$ and $f(0) = 0$, so $f = 0$; and $Vf = 0$ forces $\int_0^xf = 0$ for all $x$, hence $f = 0$ almost
everywhere.
\item The \emph{right shift} $S(x_1, x_2, \ldots) = (0, x_1, x_2, \ldots)$ on $\ell^2$ is an isometry without eigenvalues ($Sx = \lambda x$ gives $\lambda x_1 = 0$, $\lambda x_{k+1} = x_k$, hence $x = 0$), and it is not surjective. Its adjoint, the left shift, has
every $\lambda$ with $|\lambda| < 1$ as an eigenvalue, with eigenvector $(1, \lambda, \lambda^2, \ldots)$. So an injective bounded operator on an infinite-dimensional space need not be surjective, unlike in finite dimensions.
\end{enumerate}
\end{counterexample}

\section{Fr\'echet spaces}\label{functional-analysis:section-frechet}

The space $C^\infty(U)$ of smooth functions on an open set $U \subseteq \RR^n$ has a natural notion of convergence, uniform convergence of all derivatives on compact sets, which does not come from a
norm. It is described by countably many seminorms.

\begin{definition}\label{functional-analysis:definition-frechet}
Let $X$ be a vector space with a sequence of seminorms $p_1 \leq p_2 \leq \cdots$ (Theorem~\ref{functional-analysis:theorem-hahn-banach}) that \emph{separates points}: $p_j(x) = 0$ for all $j$ only if
$x = 0$. The sets $x + \{y : p_j(y) < \varepsilon\}$ form a basis of a topology, in which $x_k \to x$ if and only if $p_j(x_k - x) \to 0$ for every $j$. It is induced by the translation invariant metric
\[
d(x, y) = \sum_{j \geq 1}2^{-j}\min\bigl(1, p_j(x - y)\bigr),
\]
and $X$ is a \emph{Fr\'echet space} if this metric is complete.
\end{definition}

That the sets form a basis and that $d$ induces the topology follow from the inequalities: $d(x, y) < 2^{-j}\varepsilon$ implies $p_j(x - y) < \varepsilon$ for $\varepsilon \leq 1$, and conversely
$p_j(x - y) < \varepsilon$ implies $d(x, y) < \varepsilon + 2^{-j}$ (the terms with index $\leq j$ are at most $2^{-i}\varepsilon$ since the seminorms increase, and the tail is at most $2^{-j}$).
Banach spaces are Fr\'echet spaces with $p_j = \|\cdot\|$ for all $j$.

\begin{example}\label{functional-analysis:example-frechet}
\begin{enumerate}
\item $C(\RR)$ with $p_j(f) = \sup_{[-j, j]}|f|$: convergence is uniform convergence on compact sets, and completeness follows from that of $C([-j, j])$ for each $j$.
\item $C^\infty(U)$, with an exhaustion $K_1 \subseteq K_2 \subseteq \cdots$ of $U$ by compact sets ($K_j \subseteq \operatorname{int}K_{j+1}$, $\bigcup K_j = U$) and
$p_j(f) = \sum_{|\alpha| \leq j}\sup_{K_j}|\partial^\alpha f|$. A Cauchy sequence has all derivatives uniformly Cauchy on each $K_j$; the limits are continuous, and they are the derivatives of the limit
of the functions, by the theorem on uniform convergence of derivatives (Theorem~\ref{real-analysis:theorem-uniform-derivative}) applied in each variable. So $C^\infty(U)$ is a Fr\'echet space.
\item The \emph{Schwartz space} $\mathcal{S}(\RR^n)$ of smooth functions all of whose derivatives decay faster than every power, with $p_j(f) = \sup_x(1 + |x|)^j\sum_{|\alpha| \leq j}|\partial^\alpha f(x)|$; it
contains $e^{-|x|^2}$ and is the domain of the Fourier transform in Chapter~\ref{distributions}.
\end{enumerate}
\end{example}

\begin{proposition}\label{functional-analysis:proposition-not-normable}
The topology of $C(\RR)$ in Example~\ref{functional-analysis:example-frechet}(1) is not induced by any norm.
\end{proposition}

\begin{proof}
Suppose a norm induces it, with unit ball $W$. Then $W$ contains a basic neighbourhood $\{f : p_j(f) < \varepsilon\}$, and the neighbourhood $U = \{f : p_{j+1}(f) < 1\}$ of $0$ contains a ball $tW$,
$t > 0$, of the norm. Let $\phi$ be continuous, supported in $(j, j + 1)$, with maximum $1$. For every $c > 0$, $p_j(c\phi) = 0$, so $c\phi \in W$ and $tc\phi \in tW \subseteq U$; but
$p_{j+1}(tc\phi) = tc$, which exceeds $1$ for $c$ large. Contradiction.
\end{proof}

The Baire category theorem applies to Fr\'echet spaces, since they are complete metric spaces, and the uniform boundedness, open mapping and closed graph theorems hold for them with the same proofs
\cite[Chapter 2]{rudin-fa}. A linear map $T \colon X \to Y$ of Fr\'echet spaces, with seminorms $(p_k)$ and $(q_j)$, is continuous if and only if for every $j$ there are $k$ and $C$ with $q_j(Tx) \leq Cp_k(x)$
\cite[Chapter 1]{rudin-fa}; for example
$\partial/\partial x^i \colon C^\infty(U) \to C^\infty(U)$ is continuous, with $p_j(\partial_if) \leq p_{j+1}(f)$, although differentiation is unbounded for every norm on $C^1$
(Counterexample~\ref{functional-analysis:counterexample-unbounded}).

\section{Weak topologies}\label{functional-analysis:section-weak}

\begin{definition}\label{functional-analysis:definition-weak}
The \emph{weak topology} on a normed space $X$ is the initial topology (Definition~\ref{general-topology:definition-initial-final}) for all $\Lambda \in X^*$; $x_n \rightharpoonup x$ weakly means
$\Lambda(x_n) \to \Lambda(x)$ for all $\Lambda$. The \emph{weak-$*$ topology} on $X^*$ is the initial topology for the evaluations at points of $X$.
\end{definition}

\begin{theorem}[Banach--Alaoglu]\label{functional-analysis:theorem-alaoglu}
The closed unit ball of $X^*$ is weak-$*$ compact. If $H$ is a separable Hilbert space, every bounded sequence in $H$ has a weakly convergent subsequence.
\end{theorem}

\begin{proof}
Map $\Lambda \mapsto (\Lambda(x))_{x \in X}$ into $\prod_{x \in X}\overline{D}(0, \|x\|)$, which is compact by Tychonoff's theorem (Theorem~\ref{general-topology:theorem-tychonoff}). The image consists of the
families satisfying the (closed) linearity conditions $\lambda_{ax+by} = a\lambda_x + b\lambda_y$, so it is a closed subset, hence compact, and the map is a homeomorphism onto its image by definition of the
weak-$*$ topology. For the second statement let $(x_n)$ be bounded and $(y_k)$ a dense sequence in $H$. By a diagonal argument, after passing to a subsequence $\langle x_n, y_k \rangle$ converges for every $k$,
hence $\langle x_n, y \rangle$ converges for every $y$ (uniformly approximate $y$ by the $y_k$, using boundedness). The limit is a bounded conjugate-linear functional of $y$, so it is $\langle x, y \rangle$ for
some $x$ by Theorem~\ref{functional-analysis:theorem-riesz-representation}.
\end{proof}

\begin{counterexample}\label{functional-analysis:counterexample-weak}
Weak convergence does not imply norm convergence: an orthonormal sequence $(e_n)$ in a Hilbert space satisfies $e_n \rightharpoonup 0$ by Bessel's inequality, while $\|e_n\| = 1$. By
Proposition~\ref{functional-analysis:proposition-banach-examples}(3), the closed unit ball of an infinite-dimensional space is never norm compact, although it is weakly compact for Hilbert spaces.
\end{counterexample}

\section{Exercises}\label{functional-analysis:section-exercises}

\begin{exercise}\label{functional-analysis:exercise-adjoint-properties}
Show that $(ST)^* = T^*S^*$, $T^{**} = T$ and $\|T^*T\| = \|T\|^2$ for bounded operators on a Hilbert space.
\end{exercise}

\begin{solution}
The first two follow from uniqueness of adjoints. For the last, $\|T^*T\| \leq \|T\|^2$, and $\|Tx\|^2 = \langle T^*Tx, x \rangle \leq \|T^*T\|\,\|x\|^2$.
\end{solution}

\begin{exercise}\label{functional-analysis:exercise-integral-operator}
Let $k \in L^2([0,1]^2)$ and $(Tf)(x) = \int_0^1k(x, y)f(y)\,dy$. Show that $T$ is a compact operator on $L^2([0,1])$ with $\|T\| \leq \|k\|_2$, and that $T^*$ has kernel $\overline{k(y, x)}$.
\end{exercise}

\begin{solution}
Boundedness follows from Cauchy--Schwarz and Fubini's theorem. Kernels of the form $\sum_i\phi_i(x)\psi_i(y)$ give finite rank operators and are dense in $L^2([0,1]^2)$ (products of an orthonormal basis of
$L^2([0,1])$ form an orthonormal basis of $L^2([0,1]^2)$), so $T$ is a limit of finite rank operators. The formula for $T^*$ is Fubini's theorem.
\end{solution}

\begin{exercise}\label{functional-analysis:exercise-no-hahn-banach-norm}
Show that the functional $f \mapsto f(0)$ on $C([-1,1])$ is not continuous for the $L^2$ norm, and deduce that it does not extend to $L^2([-1,1])$.
\end{exercise}

\begin{exercise}\label{functional-analysis:exercise-spectral-projections}
Let $T$ be compact self-adjoint with eigenvalues $\lambda_k$ and orthonormal eigenvectors $e_k$. Show that $Tx = \sum_k\lambda_k\langle x, e_k \rangle e_k$, that $T$ is a norm limit of the finite rank operators
obtained by truncating this sum, and that $\|T\| = \max_k|\lambda_k|$.
\end{exercise}

\begin{exercise}\label{functional-analysis:exercise-l1-incomplete}
Show that $C([0, 1])$ with the norm $\|f\|_1 = \int_0^1|f|$ is not complete, by exhibiting a Cauchy sequence converging in $L^1$ to $\mathbf{1}_{[1/2, 1]}$.
\end{exercise}

\begin{solution}
Let $f_n$ be $0$ on $[0, \frac12 - \frac1n]$, $1$ on $[\frac12, 1]$ and linear in between. Then $\|f_n - f_m\|_1 \leq \frac1{\min(n, m)}$, and $\|f_n - \mathbf{1}_{[1/2, 1]}\|_1 \leq \frac1n$. If $f_n \to f$ in
$\|\cdot\|_1$ with $f$ continuous, then $f = \mathbf{1}_{[1/2, 1]}$ almost everywhere, which is impossible: $f$ would be $0$ on $[0, \frac12)$ and $1$ on $(\frac12, 1]$ by continuity, and could not be
continuous at $\frac12$.
\end{solution}

\begin{exercise}\label{functional-analysis:exercise-multiplication-operator}
Let $g \in L^\infty([0, 1])$. Show that $M_gf = gf$ defines a bounded operator on $L^2([0, 1])$ with $\|M_g\| = \|g\|_\infty$, that $M_g^* = M_{\bar g}$, and that $M_g$ is compact only if $g = 0$.
\end{exercise}

\begin{solution}
$\|gf\|_2 \leq \|g\|_\infty\|f\|_2$. For $\varepsilon > 0$ the set $E = \{|g| > \|g\|_\infty - \varepsilon\}$ has positive measure, and $f = \mathbf{1}_E$ gives $\|gf\|_2 \geq (\|g\|_\infty - \varepsilon)\|f\|_2$.
$\langle gf, h\rangle = \int gf\bar h = \langle f, \bar gh\rangle$. If $g \neq 0$, some $E = \{|g| > \delta\}$ has positive measure, and $L^2(E)$ is infinite-dimensional (it contains the indicator functions
of disjoint subsets of $E$ of positive measure, which exist since $t \mapsto \lambda(E \cap [0, t])$ is continuous). On $L^2(E)$, $M_g$ has the bounded inverse $M_{1/g}$, so if it were compact the
identity of $L^2(E)$ would be compact, contradicting Proposition~\ref{functional-analysis:proposition-banach-examples}(3).
\end{solution}

\begin{exercise}\label{functional-analysis:exercise-bessel-equality}
Let $(e_n)$ be an orthonormal sequence in a Hilbert space and $(c_n) \in \ell^2$. Show that $\sum_nc_ne_n$ converges, that its norm is $(\sum|c_n|^2)^{1/2}$, and that $c_n = \langle x, e_n\rangle$ for its
sum $x$ (the Riesz--Fischer theorem for orthonormal sequences).
\end{exercise}

\begin{solution}
By Pythagoras $\|\sum_{n=M}^Nc_ne_n\|^2 = \sum_{n=M}^N|c_n|^2$, so the partial sums are Cauchy and converge by completeness, with the stated norm in the limit. By continuity of the inner product,
$\langle x, e_m\rangle = \lim_N\sum_{n \leq N}c_n\langle e_n, e_m\rangle = c_m$.
\end{solution}

\begin{exercise}\label{functional-analysis:exercise-legendre}
Compute the next orthonormal polynomial $P_3$ of Example~\ref{functional-analysis:example-bases-l2}(2), and find the polynomial $q$ of degree at most $2$ minimising $\int_0^1(e^t - q(t))^2\,dt$.
\end{exercise}

\begin{solution}
Gram--Schmidt on $x^3$ gives $P_3 = \sqrt7(20x^3 - 30x^2 + 12x - 1)$: one checks that $20x^3 - 30x^2 + 12x - 1$ is orthogonal to $1$, $x$, $x^2$ (its integrals against them are $5 - 10 + 6 - 1 = 0$,
$4 - 7.5 + 4 - 0.5 = 0$ and $\frac{20}{6} - 6 + 3 - \frac13 = 0$) and has norm$^2$ $\frac17$. The best approximation is the orthogonal projection $q = \sum_{n \leq 2}\langle e^t, P_n\rangle P_n$, with
$\langle e^t, P_0\rangle = e - 1$, $\langle e^t, P_1\rangle = \sqrt3(3 - e)$ and $\langle e^t, P_2\rangle = \sqrt5(7e - 19)$, using $\int_0^1te^t = 1$ and $\int_0^1t^2e^t = e - 2$.
\end{solution}

\begin{exercise}\label{functional-analysis:exercise-frechet-polynomials}
Show that a sequence of polynomials converging in $C^\infty(\RR)$ (Example~\ref{functional-analysis:example-frechet}(2)) to a function $f$ need not have bounded degrees, but that if the degrees are bounded
by $d$, then $f$ is a polynomial of degree at most $d$. Deduce that the polynomials of degree at most $d$ form a closed subspace, and that all polynomials form a dense, non-closed subspace of $C^\infty(\RR)$.
\end{exercise}

\begin{solution}
The partial sums of $\sum x^k/k!$ converge to $e^x$ in $C^\infty(\RR)$, as all derivatives converge uniformly on compact sets. If $\deg p_k \leq d$, then $p_k^{(d+1)} = 0$ converges to $f^{(d+1)}$, so
$f^{(d+1)} = 0$ and $f$ is a polynomial of degree at most $d$ (Theorem~\ref{real-analysis:theorem-mean-value} applied repeatedly). Density: by the Weierstrass theorem applied to $f^{(j)}$ on $[-R, R]$ and
integrating $j$ times, every smooth function is approximated with its first $j$ derivatives on $[-R, R]$ by polynomials; the space is not closed since $e^x$ is not a polynomial.
\end{solution}
